<html>
<head>
<title> Section P. RUMP
</head>

<body>

<h2>Important Notes</h2>

<ul>
<li> RUMP as distributed includes stopping power tables
appropriate for incident energies to 3.0 MEV.  If you are working above this
energy, see appendix on QSTOPP and STOPP for generating higher energy tables.
<li> RUMP does not normally handle resonances since these are very
geometry dependent.  To include non-Rutherford cross sections, you must
create a {\tt .ADT} file as described in appendix on RES5.
<li> Forward scattering simulations requires manual entry of cross
sections.  See commands {\tt XSECT} and {\tt E2COEF}.
<li> The native format for {\tt .RBS} files has changed. The old
format is still recognized by RUMP.  However translator programs from
{\tt .RBS} to other formats will need to be rewritten.  See the appendix on
the new data format.
</ul>

<h2> Introduction and Acknowledgements </h2>

<p> 
RUMP is a series of FORTRAN subroutines designed for the analysis and
simulation of Rutherford Backscattering Spectroscopy (RBS) data.  It was
initially developed at Cornell University in Dr. J.W. Mayer's research group
by then graduate students L.R. Doolittle and M.O. Thompson with assistance
from R.C. Cochran.  The development of the efficient simulation algorithm and
the semi-automatic search algorithms formed a major part of L.R. Doolittle's
thesis.  Numerous other graduate students and researchers at Cornell and other
sites have since contributed to the program.  We wish to acknowledge Dr.\
E.J. Kramer's group at Cornell for introducing equations to SIM as well as
being the impetus for the forward scattering support. Dr. P. Revesz is also
acknowledged for numerous improvements in both the code and user interface
including most recently the smoothing algorithms.  Finally, we most gratefully
acknowledge all of the past and current graduate students at Cornell who have
used the program for real analysis and data reduction.  We firmly believe that
the value of RUMP stems directly from the fact that it is used by the same
people who are responsible for its development.

<p>
The algorithms and subroutines comprising RUMP are considered to be in the
public domain.  The source code for all RUMP specific routines are included in
this distribution.  They may be freely modified, recompiled or distributed
as long as the original copyright notice is retained in all source code.  CGS
will act as a general repository for code corrections and enhancements and
will attempt to disseminate modifications in an orderly manner.  However,
while the code to RUMP is public domain, the library support routines are part
of a commercial product licensed for use on one machine only and may not be
distributed.

<p> The meaning of the acronym RUMP is a question often raised in discussions.
Let us set the record straight; the name was invented solely to avoid
trademark infringement - the earliest name of this package was TRACOR.
However, for those who demand full accounting and explanations, we suggest the
following expansions:

<ul>
  <li> Rutherford Universal Manipulation Program
  <li> RBS Universal Master Package
  <li> RBS for the Upwardly Mobile Professional
  <li> Rapid Unearthing of Multiple Parameters
  <li> Random Utility for Making Pictures
  <li> Rather Unusual Masochistic Package
</ul>

<h2> Pedigree:  Relationship between RUMP and GENPLOT </h2>

<p>
RUMP is a data manipulation package similar to GENPLOT but is
designed specifically for the analysis and simulation of Rutherford
Backscattering Spectroscopy (RBS) data.  Both packages were initially
written around the same graphics and token processing libraries to
provide a consistent (but not identical) user interface.  The two
packages diverged in recent years however as GENPLOT became a
commercial product for the microcomputer environment and RUMP
development continued on a VAX mainframe.  To reduce the continuing
effort required to support separate software packages, RUMP has once
again been modified to conform to the current library support.  With
this move, RUMP support and documentation (at some meager level) has
been assumed by CGS.

<p> 
RUMP and GENPLOT are designed on equivalent software levels --- both are
the upper shell of complex data analysis and graphics packages.  RUMP is
neither a subset of GENPLOT nor is GENPLOT a subset of RUMP.  They are
equivalent subroutines utilizing a common set of graphics and system support
libraries.  Indeed, with sufficient memory GENPLOT is a command within RUMP
and RUMP can be a command within GENPLOT.  The commands available in either
packages can be divided into two classes; upper level commands implemented
within the main program itself and lower level commands implemented in the
underlying library.  {\tt PLOT} is an upper level command and hence operates
in a different manner in each package.  However, {\tt SCRIPT} and {\tt
ANNOTE} are support library commands and are identical in both RUMP and
GENPLOT.  It is important to remember this distinction if you will be using
both packages or are already familiar with one.

<p>
This section of the manual deals only with commands which are specific to
RUMP.  Lower level commands common to RUMP and GENPLOT are fully described in
other sections of the manual.  We have attempted to clearly identify which
commands operate identically between the two packages as well as point out
specific incompatibilities in the commands.  Even if you are an old user of
RUMP, some commands have changed considerably and we ask that you {\it please
read through section 1 of the GENPLOT manual} paying close attention to the
``Essential Information'' (which is repeated in the ``Quick Reference''
section).

<h2> Algorithm and Data References </h2>

<ol>
  <li> <it>Algorithms for the Rapid Simulation of Rutherford Backscattering
       Spectra</it>, L.R. Doolittle, Nucl. Inst. Meth. {\bf B9}, 344 (1985).
  <li> <it>A Semiautomatic Algorithm for Rutherford Backscattering Analysis</it>,
       L.R. Doolittle, Nucl. Inst. Meth. {\bf B15}, 227 (1986).
  <li> <it>Empirical Stopping Powers for Ions in Solids</it>, J.F. Ziegler, J.P. 
       Biersack, and U. Littmark, IBM Research Report RC9250 (1982).
  <li> <it>Numerical Recipes</it>, Press, Flannery, Teukolsky and Vetterling,
       (Cambridge University Press, 1986).
  <li> <it>Backscattering Spectrometry</it>, W.K. Chu, J.W. Mayer and M.-A.
       Nicolet, (Academic Press, NY, 1978).
  <li> <it>Ion Beam Handbook for Material Analysis</it>, J.W. Mayer and E. 
       Rimini, (Academic Press, NY, 1977).
  <li> L.R. Doolittle, Ph.D Thesis, (Cornell University, Ithaca, NY 1987).
</ol>

<h2> General Concepts: Buffers and Data Files </h2>

@@{Buffers and Files--RUMP}

RBS experimental and simulation spectra are stored in RUMP in a set of
internal buffers.  These buffers are reused on a rotating basis to accommodate
new spectral data.  Twelve buffers are configured in the PC version, each
with a maximum of 1024 channels of data (mainframe versions generally increase
the buffer size to 2048 channels).  The buffers are numbered from -1 through
10.  The first two buffers, -1 and 0, are always reserved for temporary
calculations (-1) and for spectral simulations (0).  The remaining ten buffers
hold experimental spectra, copies of simulations or temporary calculations
(PERT).  Buffers are referred to by either number or by content (the filename
of the corresponding data).  At any point in time, one buffer is designated as
the MAIN or ACTIVE buffer.  This designation is important in that many
commands, for example {\tt THICKNESS}, operate only on the ACTIVE buffer.  The
ACTIVE buffer may be changed using commands such as {\tt GET} or {\tt POINTAT}
and can be identified using the {\tt BUFFERS} command.  The first two special
buffers -1 and 0 cannot, in general, be retained as the ACTIVE buffer.

Buffers are actually data structures containing not only the spectral
information (counts per channel), but also experimental parameters required
to analyze the data.  These experimental parameters describe the incident
particle, the beam energy, charge integration, scattering geometry, etc.  This
information is stored on a per buffer basis and may be modified using commands
described in this manual.

RUMP data files correspond closely to the buffer structure.  The most recent
version of the data file utilize a format which is machine independent,
allowing direct binary exchange of data files between different mainframe and
micro implementations of RUMP.  An appendix to this section describes the
internal data file format in detail.  The old format, written using FORTRAN
unformatted I/O records, is still recognized by the READ commands.  Buffers
are filled with experimental data by reading the files using {\tt GET}, {\tt
REAL}, {\tt PLOT} or {\tt OVERLAY}.  The ACTIVE buffer can be saved, with the
current buffer experimental parameters, using {\tt REWRITE}, {\tt SAVE} or
{\tt WRITE} commands.


<h2> General Concepts: Specifying Elements </h2>

Elements names are required in many of the commands to RUMP.  Internally,
RUMP maintains a list of 92 elements with isotope masses, abundances, and
atomic density.  Elements are specified by their common one or two letter
symbol such as Si or O.  Full names are not permitted.  A specific isotope
can also be specified for most commands by using the format 28Si or 31Si.  (An
obsolete format Si+28 is also still supported.)  If the isotope is listed in
the internal tables, the actual mass for the element is used.  Otherwise, the
specified atomic weight will be used as the atomic mass.  42Si is valid in
RUMP but don't expect meaningful results.

Charged beams are specified in the same manner except for the addition of
plus signs after the symbol.  4He++ is the symbol for normal double plus
backscattering alpha particles.


<h2> General Concepts: SIM and PERT </h2>

RUMP includes a very fast and powerful, if not overly friendly, simulations
package referred to as SIM.  SIM couples with an even less friendly package
called PERT which automatically varies parameters describing the simulation to
reach a ``best fit'' with a given experimental spectrum.  Both SIM and PERT
are sub-processors under RUMP with their own set of commands and data files.
There is no command to do a simulation; simulations are automatically
performed when necessary using experimental parameters from the ACTIVE spectra
and elements and layer thicknesses from a ``sample description''.  Commands
within SIM are primarily directed to generating the sample description and
modifying some of the simulation parameters.  PERT is solely concerned with
the search procedure.

\newpage

<h2> Basic Plotting Commands </h2>

\bigskip

\plaintopline \autosizetable
{\tts PLot} & Erase screen and plot spectrum from buffer or file\cr
{\tts OVerlay} & Overlay spectrum from file or buffer on current plot\cr
{\tts REplot} & Plot the currently active spectrum \cr
{\tts AXis} & Draw axes only \cr
{\tts COMPare} & Plot active and simulation buffers \cr
{\tts BLowup} & Plot region of spectrum on an expanded vertical scale \cr
{\tts EXpand} & Set region to a subset of current region and replot \cr
{\tts BYe} & Quit and exit program completely\cr
{\tts Quit} & Quit and exit program completely\cr
\endautosizetable \medskip

<h2> Plot </h2>

\syntax{PLot [file | buffer \# | SIM | MAIN | ALT | LAST | NOW | THeory]}{Plot data file or specific buffer}

@@{PLOT--RUMP}

The {\tt PLOT} command is the basic command for drawing axes and spectrum to
the active graphics device.  Normally, the screen is erased (or a new page
obtained), the axes are drawn and the requested spectrum is plotted.  The data
are drawn using the current line type and symbol type.  If the autoids option
is selected, the identifier string from the spectrum will also be drawn on the
plot.

A filename or buffer number must be specified in RUMP version of the {\tt
PLOT} command.  If a number between 0 and 10 is specified, it is assumed to
refer to a buffer and that buffer plotted.  The simulation buffer will be
updated if necessary.  If any buffer other than 0 is plotted, it becomes the
ACTIVE buffer after this command.  If the simulation buffer 0 is plotted, the
ACTIVE buffer remains unchanged.

If a buffer number is not specified in {\tt PLOT}, the token is compared first
to a number of special names.  The names {\tt SIM}, {\tt ALT} and {\tt THeory}
refer to the simulation buffer 0 and hence {\tt PLOT THEORY} is equivalent to
the command {\tt PLOT 0} described above.  The name {\tt MAIN} likewise refers
to buffer 1 while {\tt NOW} and {\tt LAST} refer to the currently ACTIVE
<it>.  {\it These names should not be used for normal data files.</it>

If the token following {\tt PLOT} is neither a number nor a special name, it
is assumed to be a valid filename with a default extension of {\tt .RBS}.  The
current buffers are first searched for one already containing that filename
(full pathname is checked).  If found, the buffer becomes ACTIVE and {\tt
PLOT} operates as described above.  If no buffer currently contains the file,
the buffers are scrolled down to create an empty buffer 1 and the file is read
into buffer 1.  The original buffer 10 is lost in the process (even if the
file read fails).  The ACTIVE pointer is set to buffer 1 and the spectrum is
plotted.

The selection of a buffer occurs before any plotting.  The scale for the YIELD
axis (if autoranging) and the conversion from channel number to energy are
determined from the plotted buffer.  The channel number range is converted to
the equivalent energy range through the {\tt CONV} parameters {\tts SCALE1} and
{\tts SCALE2}.  For a {\tt PLOT}ed spectra, both energy and channel scales are always
correct.  

The yield can be plotted either as a raw counts per channel or as a normalized
yield, selected by the {\tt RAW} or {\tt NORMALIZED} commands.  In the
normalized mode, the yield is normalized to conditions of 1 keV channel energy
width, 1 $\mu$C of single plus equivalent charge, and a 1 msr detector solid
angle.  This normalization should allow direct comparison of any spectra
obtained at the same incident energy in the same scattering geometry.  (Note:
Charge equivalency treats 2~$\mu$C of He$^{++}$ the same as 1~$\mu$C of
He$^+$.)

\newpage
<h2> Overlay </h2>

\syntax{OVerlay [file | buffer \# | SIM | MAIN | ALT | LAST | NOW | THeory]}{Overlay a spectrum on existing axes}

@@{OVERLAY--RUMP} 

{\tt OVERLAY} operates identically to {\tt PLOT} with the exception that no
axes are drawn.  Existing energy and yield scales are used allowing
comparison of numerous spectra.  Plotting is always done based on the energy 
scale, {\bf not} the channel scale allowing overlay of data from different
multichannel analyzers.  The channel axes will correct only for the first
spectrum plotted unless all spectra have the same {\tt CONV} constants.  
The energy scale is always correct however.  Since {\tt OVERLAY} also uses the
current yield range, data may extend offscale if the range is inappropriate.

<h2> Replot </h2>

\syntax{REplot}{Replot most currently ACTIVE buffer}

@@{REPLOT--RUMP}

The {\tt REPLOT} command is synonymous with {\tt PLOT NOW}.  It quickly
replots the ACTIVE buffer with a little less typing.

<h2> Axis </h2>

\syntax{AXis}{Erases screen and draws the axes}

@@{AXIS--RUMP}

The {\tt AXIS} command normally erases the screen and displays a new
set of axes.  The conversion between energy and channel number and the yield
scale are determined from the currently ACTIVE buffer (see {\tt PLOT} above).
If the ACTIVE buffer contains no valid spectra data, a warning message
indicating that the axes values are out of range may be issued.

<h2> Compare </h2>

\syntax{COMPare}{Compare ACTIVE buffer to simulation buffer}

@@{COMPARE--RUMP}

The {\tt COMPARE} command is another command designed to reduce typing.  It is
equivalent to {\tt PLOT NOW LTYPE 1 OV THEORY LTYPE <old>}.  The ACTIVE buffer
is plotted in the current line and symbol types and the simulation buffer is
overlaid as a solid line.  The line type is returned to its previous value.


<h2> Blowup </h2>

\syntax{BLowup <cursor marks> <factor>}{Expand vertical scales for region of data}

@@{BLOWUP--RUMP}

The {\tt BLOWUP} command allows you to expand the vertical scale for a portion
of your data to see, for example, a small impurity peak.  The cursor marks a
region to be expanded and an expansion factor is read from the command line.
Data from the current ACTIVE buffer is overlaid on the plot using the current
symbol or line type over only the identified region with the specified
expansion factor.  The channel range may be selected using either the cursor
or the keyboard (press {\tt <ESC>} to escape from cursor to keyboard input
mode).  The actual data in the ACTIVE buffer is {\bf not} modified and hence
subsequent plots are not automatically expanded.

<h2> EXpand </h2>

\syntax{EXpand <cursor marks>}{Expand plot show only a subset region}

@@{EXPAND--RUMP}

{\tt EXPAND} is a useful combination of two operations: a) use of the cursor
to define a new region and  b) a replot of the ACTIVE buffer.
Note that {\tt EXPAND} redefines the plot region as in the {\tt REGION}
command below and that only the ACTIVE buffer is redrawn.  Other spectra on
the current plot will not be redrawn over the expanded range.


<h2> Quit </h2>

\syntax{QUIT or BYE}{Exits back to system from the RUMP program}

@@{QUIT--RUMP}

This command should be obvious (hopefully).  All buffer and status information
is lost when you exit, as are any modifications to buffers which have not been
saved.  If a HCOPY, temporary MACRO or SCRIPT file is active, they will be
disposed of as appropriate (possibly with user input).

\newpage

<h2> Plot Scaling and Axes Layout </h2>

\bigskip \plaintopline \autosizetable
{\tts PARMS} & Display current plot parameters\cr
{\tts REGion} &  Set the channel range\cr
{\tts FOrcex} & Force channel scale to be exactly as specified \cr
{\tts COunts} &  Set the range of the YIELD axis\cr
{\tts LInear} &  Specify linear mode for YIELD axis\cr
{\tts SQrt} &  Specify square root mode for YIELD axis\cr
{\tts LOg} &  Specify Logarithmic mode for YIELD axis\cr
{\tts NOrmaliz} &  Specify YIELD to be plotted in normalized units\cr
{\tts RAw} &  Specify YIELD to be plotted as raw data \cr
{\tts LAbels} & Set degree of labeling for axes \cr
\endautosizetable \bigskip

<h2> Parms </h2>

\syntax{Parms}{Displays some internal plotting status parameters}

@@{PARMS--RUMP}

The {\tt PARMS} commands displays some of the plotting and scaling information
to the screen.  See Section D for additional information on the plot size and
device information.   The lower section indicates the range of the channel 
data output, the scaling for the yield (Max Counts), type of yield scale
(LINEAR, SQRT or LOG), and the current linetype, symbol and npoint
values.

\medskip\beginscreen
 Plot device: AUTO
 Size:      9.00 by   7.30           Shrink:  1.000
 Offset:    0.50 by   0.00           Pen:      1 of max   7
 Margin:    0.85 by   0.85           Speed:   18
 Hard copy not initialized

  Min-Max Channels:   100.0  600.0     Max Counts:  Autoscaled
  Scale is linear, normalized yield.
  Linetype:  0   Symbol:  0   Npoints:  1

 Minor tick marks are disabled.
 Current directory: C:{\bs}RUMP{\bs}EXAMPLES
 Session SCRIPT file: r.cmd
\endscreen

<h2> Region </h2>

\syntax{REgion <lower-channel> <upper-channel>}{Set the channel range}

@@{REGION--RUMP}

Set the region of the lower axes in channel numbers.  The region to be
plotted is always specified in terms of channels but plotting is actually
performed in terms of energy.  The specified channels may be rounded off to
close values to give nicer plot scales.  See {\tt FORCEX} and {\tt SGRAPH}
to disable the roundoff and force the scales to cover a specified region.


<h2> Forcex </h2>

\syntax{FOrcex [ON | off]}{Force channel numbers to be exact}

@@{FORCEX--RUMP}

Normally, the channel numbers specified by the {\tt REGION} command or by
an {\tt EXPAND} command are ``rounded out'' to closest reasonable whole
numbers.  This roundoff feature can be disabled by setting {\tt FORCEX ON}.
Be sure though, when enabled, not to specify a zero length region.  The
extension of the lower boundary of a {\tt REGION} command to zero can also
be disabled (or controlled) by setting the ZFORCE variable using {\tt SGRAPH}.

<h2> Counts </h2>

\syntax{COunts <max-yield-range>}{Set the range for yield axis of the plot}

@@{COUNTS--RUMP}

The range for the YIELD axis is controlled by the {\tt COUNTS} commands.  The
YIELD axis will always begin at zero and extend to the specified maximum.  If
the {\tt COUNTS} is specified as zero, RUMP will perform autoranging of the
axis so the entire plotted range of the ACTIVE buffer fits on screen.  If
non-zero, the specified value is used without rounding.  While the meaning of
the numbers on the yield axis depends on the mode of the axis (LINEAR, LOG, or
SQRT), setting of {\tt COUNTS} should always be in terms of the maximum linear
count.  For example, in logarithmic mode 100,000 counts for the upper limit is
a {\tt COUNTS 10000} {\bf} {\tt COUNTS 5}.  When switching between modes, if
YIELD is not in autoscale mode the current setting of {\tt COUNTS} is
``translated'' to the the new mode.  However, {\tt COUNTS} is translated when
switching between normalized and raw mode; new values must be set
corresponding to the current setting.

<h2> Linear/Log/Sqrt </h2>

\syntax{LInear} {Set Y axis mode to linear scale}
\syntax{LOg}    {Set Y axis mode to logarithmic scale}
\syntax{SQrt}   {Set Y axis mode to square root scale (uniform statistical error)}

@@{LINEAR--RUMP} @@{LOG--RUMP} @@{SQRT--RUMP}

The {\tt LINEAR} mode is the default mode for displaying the yield
information.  In this mode, peak heights and areas appear visually correct and
it is easy to estimate compositions or concentrations.  While the {\tt LINEAR}
is most common, the {\tt SQRT} mode is statistically more accurate.  In {\tt
SQRT} mode, the value plotted is the square root of the number of counts in
each channel.  In this mode, noise due to counting statistics is independent
of the absolute yield magnitude and equal visual weight is given to features
of equal statistical significance.  The {\tt LOG} mode is seldom used in RBS
but is useful for spectra covering a wide dynamic range such as PIXE data.
The common logarithm (base 10) of the data is plotted and the lower range
limit is fixed at -2 for normalized data and 0 for raw mode data.


<h2> Raw/Normalize </h2>

\syntax{RAw}{Set Y axis to display raw channel counts}
\syntax{NOrmaliz}{Set Y axis to display normalized channel data}

@@{RAW--RUMP} @@{NORMAL--RUMP}

The YIELD axis can be displayed as either the raw number of counts per channel
or as a normalized yield where experimental differences between data are
eliminated by scaling.  The commands {\tt RAW} and {\tt NORMAL} switch between
these two modes.  In raw mode, the channel data is plotted exactly as recorded
by the Multi--Channel Analyzer (MCA).  In normalized mode, the data is plotted
based on the normalization of 

$$ Yield = {{Counts * Corr} \over {Q * d\Omega * \Delta E_{chan}}} $$

The units for normalization are beam dose Q in $\mu$Coulombs equivalents of
$He^{+}$ (measured charge for $He^{++}$ is divided by two), solid angles
$d\Omega$ in milli--steradians, and the MCA channel width $\Delta E_{chan}$ in
keV.  A fudge factor {\tt CORR} allows for fine tuning of the beam charge to
compensate for inaccuracies during charge integration caused by inadequate
secondary electron suppression.

In the normalized mode, yield curves should always be directly comparable
between experiments if the beam energy and scattering geometry are the same.

<h2> Labels </h2>

\syntax{LABELS [ON | BRIEF | OFF]}{Set axis labeling mode}

@@{Labels--RUMP}

The {\tt LABELS} command sets the mode of axis labeling.  On slow plotting
terminals, its use can make eliminate the time wasted putting up the
coordinate labels on the screen.  Alternatively, the labels can be completely
suppressed in a final output allowing a draftsperson (non--sexist) to label
the plot for publication.  The default mode is on.

The three modes are:

\medskip \plaintopline \autowraptable{1in}{5in}*
ON  &    Full labels: the complete title of each axis is given,
and each major tick mark is labeled.\cr
BRIEF &  Partial labels: an abbreviated title of each axis is
given, and the two end major tick marks are labeled.
This mode is recommended for routine use on slow graphics devices as it is
faster than the ON mode. \cr
OFF  &   No labels:  No text at all is drawn on the plot.  This can
be used if a draftsman is to complete the artwork.\cr
\endautowraptable

In both the {\tt ON} and {\tt BRIEF} modes, labeling conforms to AIP
standards for publication quality plots.  Minor tick marks are normally
suppressed in publication plots and this is the default mode for RUMP.
Minor tick marks may be re--enabled using the {\tt SUBTICKS} command.


\newpage

<h2> Controlling Plot Layout and Linetypes </h2>

Most of the commands which control the plot layout and curve description 
(line type or symbol type for curves) are essentially identical between
GENPLOT and RUMP.  The following table refers to other sections of this manual
where the corresponding commands are described.  Commands marked with \dag\ 
are implemented by the graphics support libraries and consequently
are absolutely identical to the GENPLOT equivalent.  Other commands should
behave in the same manner as the GENPLOT documentation but are implemented by
RUMP instead and may consequently show subtle differences.

\medskip \plaintopline \autowraptable{.25\hsize}{.75\hsize}*
\multicol{2}{l}{\bf Section A --- Device Control }\cr
{\tts CLS\dag} & Clear the text screen\cr
{\tts DEvice\dag} & Select plotting device \cr
{\tts SPeed\dag} & Set the plotter pen speed \cr
{\tts HCopy\dag} & Save plot vectors -- allows replay on hard copy device\cr
\cr
\multicol{2}{l}{\bf Section B --- Basic Plotting}\cr
{\tts AUtoids} & Automatically use IDENT string as a legend\cr
{\tts AUTOErase\dag} & Set whether screen erased with {\tt PLOT}\cr
{\tts ERase\dag} & Erase the screen or start a new page\cr
{\tts IDs} & Put the IDENT string of buffer in the legend\cr
{\tts NPoints} & Point plotting interval\cr
{\tts TIMEStamp\dag} & Put the date on the plot\cr
\cr
\multicol{2}{l}{\bf Section D --- Line Control}\cr
{\tts COLor\dag} & Select the line color\cr
{\tts LType} & Select line type for subsequent drawing\cr
{\tts LINEWidth\dag} & Set line thickness on appropriate plotters\cr
{\tts PEn\dag}  & Select the pen color\cr
{\tts SYmbol} & Select symbol used for plotting \cr
{\tts VECsize\dag} & Set the length of dashes in dashed lines\cr
{\tts VISibili\dag} & Set line visibility\cr
\cr
\multicol{2}{l}{\bf Section E --- Plot Appearance}\cr
{\tts CHARActer} & Select character set\cr
{\tts CHRset\dag} & Select character set\cr
{\tts SUBTIcks} & Toggle subtick mark display\cr
\cr
\multicol{2}{l}{\bf Section I --- Labeling}\cr
{\tts ANNOTE\dag} & Sub--process for annotating the plot\cr
\cr
\multicol{2}{l}{\bf Section F --- Size and Position}\cr
{\tts OFfset\dag} & Controls plot offset\cr
{\tts SHrink\dag} & Scale entire plot to shrink\cr
{\tts SIZe\dag} & Set the paper size\cr
{\tts MARGins\dag} & Set the margins around the plot for labels\cr
{\tts FLIPXY\dag} & Flip the X and Y axes on the page\cr
{\tts SUB\_Plot\dag} & Adjust size/offset to create a plot within a plot\cr
\cr
\multicol{2}{l}{\bf Section M --- Miscellaneous}\cr
{\tts SGRAph\dag} & Set various plot parameters\cr
\cr
\multicol{2}{l}{\dag \ Identical to GENPLOT command}\cr
\endautowraptable


\newpage

<h2> RUMP Buffers </h2>

\bigskip \plaintopline \autosizetable
{\tts BUFfers} & Display current buffer status \cr
{\tts REad} & Point at buffer with file or read RBS file into buffer 1\cr
{\tts GET} & Point at buffer or read an RBS file into buffers\cr
{\tts RERead} & Force rereads current buffer into buffer 1 \cr
{\tts SWALLOW} & Load ASCII data into buffer as counts per channel\cr
{\tts POintat} & Point at a specific buffer by number\cr
{\tts WRITE} & Writes current buffer to another file\cr
{\tts REWrite} & Rewrites current buffer to a file with corrections\cr
{\tts SAVE} & Rewrites current buffer to a file with corrections\cr
{\tts WRASCII} & Writes current buffer to an text file\cr
{\tts EMpty} & Scrolls the data and leaves an empty buffer 1\cr
{\tts RECALcul} & Mark sample structure as modified to force recalculation\cr
{\tts RELEASE} & Release the ACTIVE buffer and scroll remaining up\cr
{\tts NEWall} & Clear all buffers and zero all data \cr
\endautosizetable \bigskip

<h2> Buffers </h2>

\syntax{BUffers}{Display RUMP buffer status}

@@{BUFFERS} @@{Buffers} @@{Buffers--RUMP}

Every {\tt PLOT}, {\tt OVERLAY}, or {\tt READ} command either points to a
spectrum in memory or reads a new file into the temporary buffer space (see
General Concepts at beginning of RUMP manual).
Buffers are referred to by number and the contents of the buffers can be
determined using the {\tt BUFFERS} command.  The {\tt BUFFERS} command
summarizes the current contents of each buffer.  Special buffers and any
unused buffers are not listed by this command.

The {\tt BUFFERS} command lists all non-empty files by buffer number.
Following the buffer number are two special flags described below.  The last
20 characters of the entire pathname for the source file are shown next
followed by the identifier text for each buffer.  The most recent buffer
plotted or otherwise acted upon is called the ACTIVE buffer.  It is designated
by the `*' in the {\tt BUFFERS} display before the buffer number.  You can
explicitly make any buffer the ACTIVE buffer using the {\tt POINTAT} command.
RUMP never points to the simulation buffer (0) unless you explicitly request
so with the {\tt POINTAT} command. 

Two other special markings used in the buffers listing indicate
modifications to the buffer.  The `\@' marking indicates files where
the experimental parameters in the buffer have been modified using
commands such as {\tt MEV} or {\tt CONV}.  Buffers marked with a
`?' indicate that the actual spectrum data (counts per channel) in the
buffer has been changed by a command such as the background
subtraction or the data smoothing.  All {\tt WRITE} commands will
prompt for verification before writing out buffers with modified
data.  No verification is necessary if only header information has
been modified.

\medskip\beginscreen
At your service! \underbar{buf}
* => Main    ? => Dirty    @ => Parameters changed
  Main    1    c:{\bs}psmith{\bs}new{\bs}wb.rbs  sput w, 300 sec in ar
  Buffer  2 ?  c:{\bs}psmith{\bs}new{\bs}wa.rbs  sput w, 300 sec in n2
 *Buffer  3    psmith{\bs}new{\bs}sputw.rbs  sputtered w, as deposited
  Buffer  4  \@ psmith{\bs}new{\bs}qq\_40.rbs  WSi on Si in Ar for 5 minutes
  Buffer  5    psmith{\bs}new{\bs}ii\_40.rbs  WSi on Si in N2
  Buffer  6    ith{\bs}new{\bs}e1-w1-n2.rbs  sample e1-w1, 300 sec in N2
Your wish? 
\endscreen


<h2> Read </h2>

\syntax{REad <filename>}{Read a spectra file from disk}

Synonym:  {\tt GET}

@@{Read--RUMP}

This command obtains a spectrum, from disk if necessary, and makes it the
ACTIVE buffer.  Either a true filename, a buffer number, or one of the special
names listed under the {\tt PLOT} command may be specified as the filename. 
If the filename already exists as one of the buffers, that buffer is simply
made active.  Otherwise, {\tt READ} scrolls the buffers creating a empty
buffer 1 and reads the file into buffer 1.  Buffer 10 is lost in the process.
The ACTIVE buffer is set to the buffer containing the file.  Each time a file
is read, the filename and the ID of the file are echoed to the screen.

Wildcard filenames may also be specified to {\tt READ}.  In the case of
wildcards, the process is repeated for every file which matches the wildcard
name.  All RBS files could be read, for example, using the command {\tt READ
*.RBS}.  Each matching file is retrieved into a unique buffers.  If more
matches occur than there are buffers, RUMP will stop reading once all the
buffers have been filled by the command.  A warning message indicating that 
the maximum number have been read is displayed in this event.

{\tt GET} is the normal manner in which the ACTIVE buffer is changed. 
Specifying {\tt GET 3} or {\tt GET aa\_40.rbs} will change the ACTIVE pointer
to the appropriate buffer.

\medskip\beginscreen
I want a cookie!! \underbar{/* Have an Oreo!}
I want a cookie!! \underbar{read {\bs}psmith{\bs}new{\bs}*.rbs}
File: c:{\bs}psmith{\bs}new{\bs}aa\_40.rbs
Id:   WSi on Si, control
File: c:{\bs}psmith{\bs}new{\bs}e-beam-i.rbs
Id:   e-beam w, as implanted
   .........
File: c:{\bs}psmith{\bs}new{\bs}wb.rbs
Id:   sput w, 300 sec in ar
WARNING: Maximum number read - All buffers filled
Your wish? \underbar{get aa\_40}
Yes master?
\endscreen


<h2> Reread </h2>

\syntax{RERead <filename>}{Read a file that was already read}

@@{REREAD--RUMP}

Under normal conditions, RUMP will always return to the buffer already
containing a given file if it is requested again.  Under some circumstances,
such as after a failed background subtract, it is necessary to bypass this and
force RUMP to read the file off the disk.  At other times, one may just want
to return to the original settings of the experimental parameters. 
{\tt REREAD} performs this action by changing the extension of the filename
associated with the ACTIVE buffer to {\tt .OLD}.  The original filename is
then reread from disk into buffer 1.  Note that the original buffer remains
but has a new name.  The original buffer 10 is lost in the shuffle since a new
file is added to the buffer stack.  It is difficult to reread a file if you
originally gave the data file an {\tt .OLD} extension.  The ACTIVE buffer
points to the reread copy of the file after this command.

\medskip\beginscreen
Next command: \underbar{reread    /* ACTIVE was buffer 2}
File: c:{\bs}psmith{\bs}new{\bs}wb.rbs
Id:   sput w, 300 sec in ar
Your wish? \underbar{buf}
* => Main    ? => Dirty    @ => Parameters changed
 *Main    1    c:{\bs}psmith{\bs}new{\bs}wb.rbs  sput w, 300 sec in ar
  Buffer  3    c:{\bs}psmith{\bs}new{\bs}wa.rbs  sput w, 300 sec in n2
  Buffer  2    c:{\bs}psmith{\bs}new{\bs}wb.old  sput w, 300 sec in ar
  Buffer  4    psmith{\bs}new{\bs}sputw.rbs  sputtered w, as deposited
You called? 
\endscreen

<h2> Swallow </h2>

\syntax{SWALLOW <channel data>}{Load data as counts of the active buffer}

@@{SWALLOW--RUMP} @@{WRASCII--RUMP} @@{ASCII data--RUMP}

The {\tt SWALLOW} command is primarily intended for macro files to load in
non-RUMP format stored ASCII channel data.  It permits the reading of ASCII
data directly into the counts array of the ACTIVE buffer.  Channel data is
read from the command line (or the active macro) one element at a time until a
blank entry (or the default value) is entered.  Multiple channels of data are
allowed on each line.  The {\tt WRASCII} command, which generates a macro file
to reload data in a pure ASCII format, uses the {\tt SWALLOW} command.  The
following macro file will load in a spectrum segment using the {\tt SWALLOW}
command.  (Note the blank line used to terminate the swallow read.)  The
command {\tt XEQ <file>} would be used to actually load the segment.

\medskip\beginscreen
/* Open a blank buffer and start swallowing the data
empty swallow
   0    0    0    0    0    0    0    0    0    0    0    0
   0    0    0    0    0    2   29  223  485  786 1049 1231
1373 1416 1504 1557 1567 1517 1539 1507 1553 1463 1516 1458
1411 1478 1326 1417 1428 1375 1344 1306 1279 1238 1219 1289
1246 1264 1192 1220 1212 1200 1166 1123 1142 1181 1137 1131
1061 1065 1076 1042 1088 1110 1070 1086 1078 1017 1057 1062

/* After a blank line to terminate SWALLOW, set parameters
mev  3.02 charge 10 covers 4.95 1.6 theta 7 phi 9 omega 3.4
corr 1.0  fwhm 25 ident 'This is test swallow data'
\endscreen

See also: {\tt WRASCII}

<h2> {REWrite </h2>

\syntax{REWrite}{Rewrites current buffer to a file with corrections}

Synonym:  {\tt SAVE}

@@{REWRITE--RUMP} @@{SAVE--RUMP}

Rewrites current buffer to the current file with corrections.
See Parameter Modifcation.

<h2> {Write </h2>

\syntax{WRITE <filename>}{Writes current buffer to another file}

@@{WRITE--RUMP}

Writes the current buffer to another file.

<h2> {WRASCII </h2>

\syntax{WRASCII <filename>}{Writes current buffer to a text file}

@@{WRASCII--RUMP}

Writes current buffer to a text file suitable for reading with {\tt SWALLOW}.

See also: {\tt SWALLOW}

<h2> Pointat </h2>

\syntax{POintat <n>}{Set the active buffer}

@@{POINT--RUMP} @@{Active Buffer--RUMP}

{\tt POINTAT} allows the ACTIVE pointer to be set to any buffer, including the
simulation and temporary buffers.  Other than the requirement that it must
point to a configured buffer, no check is made to determine if the buffer is
filled with valid data.  {\tt POINTAT} is the only command which will leave
the ACTIVE buffer pointing to the simulation or temporary buffer space.

<h2> Empty </h2>

\syntax{EMpty}{Scroll buffers to leave an empty buffer 1}

@@{EMPTY--RUMP}

{\tt EMPTY } causes the current list of buffers to be scrolled with an empty
buffer left as buffer 1.  Buffer 10 is lost in the process.  All parameters in
buffer 1 are set to defaults, the data is set to all zero, the number of
points is to the maximum and the buffer is marked as ACTIVE.  This command
can be used with the {\tt SWALLOW} command below to read in ASCII data.

<h2> RECALCULATE </h2>

\syntax{RECALculate}{Mark sample structure as modified so recalculated}

@@{RECALCULATE--RUMP}

Normally RUMP tracks the status of the theory buffer.  It marks the buffer
as being out of date if either the sample structure is modified or if the
experimental parameters are changed (such as changing the ACTIVE buffer).
However, it does not track modifications to the actual theory buffer data
(such as a {\tt BACKGROUND} subtract) nor can it track modifications to sample
parameters which are linked to the function evaluator (such as in {\tt
EQUATION USER}).  The {\tt RECALC} command is a manual override to the
automatic feature and simply marks the theory buffer as being out of date.
The next reference to the theory buffer following this command will cause the
simulation to be recalculated.


<h2> {Release </h2>

\syntax{RELease}{Release the ACTIVE buffer and unscroll}

@@{RELEASE--RUMP}

Buffers can be intentionally released for general use by the {\tt RELEASE}
command.  The data in the buffer is forgotten and all parameters for the
buffer are reset to default values.  The remaining buffers are unscrolled
to fill in the space vacated by the released buffer.  Releasing buffer 2
causes buffer 3 to become 2, 4 to become 3 etc.  It is not possible to
retrieve the data in the buffer once it is released.  You must read the file
back into memory directly.  {\tt RELEASE} is primarily used in macro files
to maintain control over which data files remain active.

<h2> NEWALL </h2>

\syntax{NEWall}{Reset all parameters and clear all buffers}

@@{NEWALL--RUMP}

The {\tt NEWALL} command ``forgets'' completely about the current buffers and
resets all parameters in the buffers to default values.  It returns the
buffer stack to the same state as when RUMP was initiated.  This command is
commonly used with macros to clear all previous reads when it is necessary to
ensure the order of files in the buffers.


\newpage

<h2> Listing and Modifying Spectrum Parameters </h2>

\bigskip \plaintopline \autosizetable
{\tts ACtive} & List the parameters of the active buffer\cr
{\tts BEAM} & Type of incident beam (He$^{++}$) \cr
{\tts CHarge} & Measured beam dose in $\mu$C \cr
{\tts CHOFF} & Channel number of first data point\cr
{\tts CONVersi} & Channel to energy conversion constants\cr
{\tts CORrecti} & Minor fudge factor for normalization correction\cr
{\tts CURRent} & Average beam current (pileup calculations)\cr
{\tts DATE} & Date of data collection\cr
{\tts FILE} & Spectrum file name\cr
{\tts FWHM} & Detector resolution (keV) \cr
{\tts GEOMetry} & Scattering geometry used \cr
{\tts IDEntifi} & Description identifier of the spectrum\cr
{\tts MEV} & Beam energy\cr
{\tts OMEGA} & Detector solid angle\cr
{\tts PHI} & Supplement of the scattering angle\cr
{\tts PSI} & Angle of surface normal to detected beam \cr
{\tts THEta} & Sample tilt or angle of surface normal to incident beam \cr
\endautosizetable \bigskip

<h2> Active </h2>

\syntax{ACtive}{List the parameters of the active buffer}

@@{ACTIVE--RUMP}  @@{Active Buffer--RUMP}

Each buffer includes information on the experimental parameters used in the
data collection process.  These parameters are also stored in the data files.
The values for the parameters of the ACTIVE buffer are listed to the screen
using the {\tt ACTIVE} command.  This listing is organized so that related
parameters are listed together on the same line.  The commands described
below allow most of the parameters to be modified.

\medskip\beginscreen
Ready if you are! \underbar{buf}
* => Main    ? => Dirty    @ => Parameters changed
  Theory  0    temp.rbs              Simulation of Ni/Ni-Si/Si
 *Main    1    rump{\bs}tests{\bs}examp.rbs  Ni/NiSi/Si Annealed 90 min
Your wish? \underbar{active}
 Filename:    c:{\bs}rump{\bs}tests{\bs}examp.rbs
 Identifier:  Ni/NiSi/Si Annealed 90 min
 LTCT Text:   LT= 857 CT= 860
 Date:        18-JUN-1985 12:33:48.48
 Beam:        3.020 MeV     4He++   10.00 uCoul   @  8.0 nA
 Geometry:    Cornell  Theta:    7.0  Phi:    9.0  Psi:    0.0
 MCA:         Econv:  4.950  1.600    First chan: 0.0   NPT: 1024
 Detector:    FWHM: 35.0 keV   Omega: 3.400
 Correction:  1.0041
Your wish? 
\endscreen

<h2> Parameter Modification </h2>

@@{FILE--RUMP} @@{IDENT--RUMP}  @@{DATE--RUMP} @@{MEV--RUMP}  @@{CHARGE--RUMP}
@@{BEAM--RUMP}  @@{CORRECTION--RUMP} @@{GEOMETRY--RUMP} @@{PHI--RUMP}
@@{THETA--RUMP} @@{PSI--RUMP}  @@{OMEGA--RUMP}   @@{CONVERSION--RUMP}
@@{CHOFF--RUMP} @@{FWHM--RUMP} @@{CURRENT--RUMP}

The following commands simply change the corresponding parameter in the
ACTIVE buffer.  The use of each parameter within RUMP is listed in the notes
below the table.  In general, only commands which one should normally modify
during a session have abbreviations.  Commands which are, or should be, solid
experimental parameters require spelling out in full - this is to avoid later
command name conflicts.

\medskip \autosizetable
Command  &   Parameter  &    Notes &   Explanation\cr
\cr

{\tt FILE}     & filename       &  -    & File from whence it came\cr
{\tt IDEntifi} & 'line of text' &  4    & Description of buffer contents\cr
{\tt DATE}     & 'line of text' &  -    & Date of data collection\cr
{\tt MEV}      & MeV            & 2,3,7 & Incident beam energy\cr
{\tt CHarge}   & $\mu$Coulombs  &  1    & Integrated beam dose\cr
{\tt BEAM}     & 4He++          & 1,2,3 & Beam type (arbitrary) \cr
{\tt CORrecti} & Value          &  1    & Correction for beam dose\cr
{\tt GEOMetry} & {\tts CORNELL/IBM/GENERAL} & 2,3,6 & Scattering geometry\cr
{\tt PHI}      & degrees        & 2,3,6 & Supplement of scattering angle\cr
{\tt THEta}    & degrees        & 2,3,6 & Sample tilt angle (incident beam)\cr
{\tt PSI}      & degrees        & 2,3,6 & Sample exit angle (exit beam)\cr
{\tt OMEGA}    & msr            &   1   & Detector solid angle\cr
{\tt CONVersi} & keV/ch \quad keV(0)  &  1,5  & ADC channel to energy calibration\cr
{\tt CHOFF}    & channel        &   5   & Channel \# of first data point\cr
{\tt FWHM}     & keV            &   3   & Detector resolution\cr
{\tt CURRent}  & nanoamps       &   3   & Beam current (for pileup)\cr
\endautosizetable 

\bigskip
{\parindent=35pt\parskip=0pt
\item{1.}  Used to normalize yield data
\item{2.}  Parameters used by RBS analytical tools
\item{3.}  Parameter used primarily by the simulation programs
\item{4.}  The {\tt IDS} command displays this text on the plot
\item{5.}  The energy scale is determined by $chan*(kev/ch) + kev(0)$.
           {\tt CHOFF} is channel number of the first data point and is
           normally set to 0. Changing {\tt CHOFF} shifts plot horizontally.
\item{6.}  Scattering geometry --- see description below.
\item{7.}  Incident beam is specified as a specific isotope and charge state
           such as 4He++.  Stopping power data for the beam should have
           previously been loaded.
\par}
 
These commands change only parameters in the active buffer.  The file remains
unchanged until a {\tt REWRITE} command updates the information in the file.
The {\tt REREAD} command may be used to restore the buffer to the state
contained in the file.

<h2> Scattering Geometry </h2>

@@{SCATTERING GEOMETRY--RUMP}

A full description of the scattering geometry for RBS requires three angles,
the scattering angle of the particle, the entrance between the surface normal
and the incident beam, and the exit angles of the beam with respect to the
surface normal.  In the geometry {\tt GENERAL}, these three angles are
independent and are set by the commands {\tt PHI}, {\tt THETA} and {\tt PSI}
respectively.  Note, however, that {\tt PHI} is really the supplement of the
scattering angle --- direct backscattering is 0 degrees.

In many cases, the three angles are not completely independent.  Given a
single tilt axis on most sample holders, only two of the angles need to be
specified and the third can be determined from knowledge of the geometry.  Two
``standard'' geometries are defined, {\tt CORNELL} and {\tt IBM}.  These are
purely historical designations having little to do anymore with the actual
geometries at the two sites.  The angle {\tt PSI}, again historically, is the
derived quantity from {\tt THETA} and {\tt PHI}.   The {\tt IBM} geometry
refers to the scattering configuration where the incident beam, surface
normal, and detected beam are all coplanar.  The exit angle {\tt PSI} is set
to {\tt PHI+THETA}.  In the {\tt CORNELL} geometry, the tilt axis lies in the
plane defined by the incident and scattered beams, and the surface normal
sweeps out a perpendicular plane.  The {\tt CORNELL} geometry sets 
$cos(psi)=cos(theta)*cos(phi)$.  If there are other common geometries in the
world, there is only space for another 32,765 possibilities in the code ---
get your request in early.

If the scattering geometry is set to {\tt GENERAL}, all three angles must
be defined and are used.  Changing the geometry to {\tt CORNELL} or {\tt IBM}
will cause {\tt PSI} to be ignored and its value may change.  

\newpage

<h2> Manipulating and Massaging Data and Buffers </h2>

\bigskip \plaintopline \autosizetable
{\tts ADD} & Add two spectra\cr
{\tts BACKground} & Background fit and subtraction\cr
{\tts COPY} & Copy one buffer to another\cr
{\tts COMPRess} & Compress channels to improve statistics\cr
{\tts MOVE} & Exchange the data in two buffers\cr
{\tts DIVide} & Calculate a $\chi_{min}$ plot\cr
{\tts SMOoth} & Smooth the data\cr
{\tts SUBTRact} & Subtract two spectra\cr
\endautosizetable
\bigskip

<h2> Move </h2>

\syntax{MOve <source buffer \#> <dest buffer \#>}{Exchange the data in two buffers}

@@{MOVE--RUMP}

The {\tt MOVE} command is really an exchange command.  The source and
destination buffers are exchanged, including all data and experimental
parameters associated with the buffers.  The ACTIVE buffer remains pointing at
the same data after this command; if either the source or destination buffer
was ACTIVE, the actual buffer number pointed to by ACTIVE will change.
The buffers {\it must} be specified by number for the {\tt MOVE} command. 
In this command only, filenames are not allowed but buffer -1 is accessible.

<h2> Copy </h2>

\syntax{COpy <source buffer> <dest buffer>}{Copy the data in source buffer to destination}

@@{COPY--RUMP}

The {\tt COPY} command directly copies all experimental parameters and data
from the source buffer to the destination buffer.  After this command, both
buffers are absolutely identical including name and modification history.  The
original contents of the destination buffer are lost.


<h2> Add </h2>

\syntax{ADd <source buffer> <dest buffer>}{Add two spectra}

@@{ADD--RUMP}

The {\tt ADD} command works with a pair of buffers, adding the contents of
the source buffer to the destination buffer and saving the result in the
destination buffer.  The source and destination buffers may be identified
either by number or by the associated filename.  {\tt ADD} will attempt to
read the file into a buffer if it is not already in the buffer set. 
However, because of the buffers scrolling, specifying a destination file
which does not already exist in the buffer set may cause the wrong buffers to
be added.  It is suggested that both files be read into buffers before
executing the {\tt ADD} command.

The result of the add is placed in the destination buffer.  This buffer will
be marked as dirty to indicate that the data has been modified (? will be
shown by the {\tt BUFFERS} command).  The {\tt ADD} command works slightly 
different depending on the setting of the {\tt RAW/NORMALIZE} switch.  In
{\tt RAW} mode the add is performed literally --- counts are counts, no
questions asked.  The charge of the destination buffer is set to the sum of
the charges of the source and and destination buffers. 

In {\tt NORMALIZED} mode, the addition is done as it appears visually on the
graph.  The charge and correction factor of the destination buffer are not
changed.  The data is summed for the same energy position (not same channel)
with corrections for the width for the channels and for the beam current,
detector and correction factors.  The resulting data is mapped onto the MCA
structure of the original destination buffer, retaining the {\tt CONV}
parameters.

\newpage
<h2> Subtract </h2>

\syntax{SUBTRact <source buffer> <dest buffer>}{Subtract two spectra}

@@{SUBTRACT--RUMP}

The {\tt SUBTRACT} command subtracts the data in the source buffer from the
destination buffer and places the results in the destination buffer.  It is
otherwise identical to the {\tt ADD} command described above and the same
comments apply.


<h2> Divide </h2>

\syntax{DIVide <source buffer> <dest buffer>}{Divide two spectra}

@@{DIVIDE--RUMP}

In the {\tt DIVIDE} command, the contents of the destination buffer are
divided channel by channel by contents of the source buffer.  Results are
stored in the destination buffer which is marked as having modified data.
See notes in {\tt ADD} for warnings on specifying filenames as buffers.
This command is commonly used to compute the ratio of a channeled spectrum to
the reference random spectrum and determine the minimum channeling yield
$\chi_{min}$.  

{\tt DIVIDE} assumes that channel to energy conversion for the two spectra are
identical, or identical enough to be ignored.  No check of the parameters is
made.  The data are divided on a channel by channel basis after correcting for
the starting channel number of each buffer.  Anything divided by zero is
defined, as far as this command is concerned, to be zero.  Only the region of
overlap between the two buffers is computed.  In {\tt RAW} mode, the absolute
channel data are divided while the {\tt NORMALIZE} mode uses the scaled yield
values.  The correction factor, solid angle and charge of the destination
buffer are all set to unity following this command.

Independent of how the mode in which the division was performed, the result
should always be plotted in {\tt RAW} mode since the normalized plots will be
completely meaningless.

<h2> Background </h2>

\syntax{BACKground <cursor entries> <order\_fit>}{Background subtract using a
polynominal fit over region}

@@{BACKGROUND--RUMP}

The {\tt BACKGROUND} command attempts to extract an isolated peak from a large
background by subtracting an estimate of the background from the data.  Peaks
extracted in this manner are then often analyzed using the {\tt THICKNESS} or
{\tt INTEGRAL} commands.  The {\tt BACKGROUND} is estimated by fitting a
polynomial, of order {\tt <order\_fit>}, through the spectral data on either
side of the isolated peak.  (You must be able to identify regions from the
background on both sides of the peak to use this command.)  Regions are
normally entered using the cursor, but may also be entered from the keyboard
by pressing the {\tt <esc>} key at the first cursor.  Four channel values are
requested and the background fit will be done in the intervals between the
lower two and upper two values.

A least squares polynomial fit is made to the data in these regions.  The
order of the polynomial must not exceed 6.  (Generally, anything above 2 is
lying to yourself --- if the data is not quadratic, don't try background
subtract.)   The polynomial fit is subtracted from the channel data over the
range of the fit, including the intervening channels.   The polynomial fit
and the resulting spectra are then overlaid on the plot.  {\it This command
modifies the data}, so one normally copies the buffer of interest into another
buffer before using the {\tt BACKGROUND} command.

The subtracted data is primarily suited to integrating.  After a background
subtract, the ``gross'' integral results will give more reliable answers than
the primitive background subtract used in the ``net'' integration.

\newpage

<h2> Compress </h2>

\syntax{COMPRess <nchan>}{Compress multiple channels to improve statistics}

@@{COMPRESS--RUMP}

The statistical noise in a spectrum can be reduced by summing adjacent
channels into a single channel using the {\tt COMPRESS} command.  The
resulting spectra has fewer data points, but each channel is statistically
less noisy.  The alternative is to collect data for a longer period initially.
The parameter {\tt <nchan>} is the number of channels which are to be combined
in each step.  Given an {\tt <nchan>} of 4, a spectrum with 1024 points will
become a spectrum with 256 points.  The energy scales and the first channel
number are modified to correspond to the compressed spectra.  The buffer is
also marked as dirty to indicate modified data.

<h2> Smooth </h2>

\syntax{SMOoth [-SV | -CONV <ntimes> | -FFT <width>]}{Smooth the data}

@@{SMOOTH--RUMP}

The smooths command reduces the noise in a spectra while attempting to
retain statistically valid features.  There are several options available
for smoothing:

{\tt \medskip \plaintopline \autosizetable
SMOOTH               &  [-RANGE <cursor or entry>]\cr
SMOOTH -SV           &  [-RANGE <cursor or entry>]\cr
SMOOTH -CONV <ntimes>& \cr
SMOOTH -FFT  <fwhm>  &  [-RANGE <cursor or entry>]\cr
\endautosizetable}

The three smoothing algorithms currently implemented are listed above, the
Savitsy-Goulay smooth ({\tt -SV}), a convolution algorithm ({\tt -CONV}) and
an FFT smoothing algorithm ({\tt -FFT}).  If no algorithm is selected, the
Savitsy-Goulay algorithm will be used by default.  The buffer is marked as
dirty to indicate modified data by this command.

The Savitsky--Goulay smooth is a simple 5 point smoothing algorithm.  It works
quickly and does a surprisingly good job.  It may be necessary to smooth
several times to achieve the desired results.  Be sure to compare the smoothed
spectra with the original before believing any results.  The exact formula is
$$
    x_n = {{-3*x_{n-2}+12*x_{n-1}+17*x_n+12*x_{n+1}-3*x_{n+2}} \over {35}}
$$

The convolution algorithm attempts to retain edges in the spectra while
smoothing.  This is a iterative algorithm which convolutes and deconvolutes
using the detector resolution on each iteration.  It will be executed the
number of times specified by the {\tt <ntimes>} value.  Typical spectra
require 5 iterations.  Note that two smooths by this algorithm are not
equivalent to one operation with twice the number of iterations.

The FFT based smoothing algorithm is based loosely on the algorithm 
in {\it Numerical Recipes}, Press et al., (Cambridge University Press, 1986).
The degree of smoothing is controlled by the specified {\tt <width>}.  The
width is roughly the number of channels over which the data will be shared
and should be approximately equal to the detector resolution divided by the
energy width of each channel.

Both the FFT and SV algorithms are implemented in a manner where the range of
smoothing can be limited to a subset of the entire buffer; the convolution
however, always operates over the entire spectra.  The {\tt -RANGE} option
specifies a finite range, marked with cursors, over which the specified smooth
will be performed.  Data outside the range will not be modified.  This feature
allows different regions of the spectra to be smoothed to differing degrees.

{\bf PERT Warning:} {\it NEVER} smooth data that PERT is to look at.
The smooth functions give pretty--looking lines for simplified viewing but
are less than useless for quantitative work.

P.S. If you are finding smooth extremely useful, you really should consider
collecting data for a longer period of time.  {\tt SMOOTH} should be used
to direct further experiments, but results should not be based on smoothed
spectra.  The {\tt COMPRESS} command above is a statistically honest method of
improving the statistics in a spectrum.

<h2> RBS ANALYSIS </h2>

\bigskip \plaintopline \autosizetable
{\tts CALibrate} & Calibrate {\tt CONV} and {\tt MEV} from spectra \cr
{\tts CURsor}  & Displays a cursor and returns energy/yield values\cr
{\tts ELement} & Indicates scattering energy for an element\cr
{\tts INFo} & Prints information about an element\cr
{\tts INTegral} & Finds integral over a region\cr
{\tts MATrix} & Plots energy and yield expected for a pure element\cr
{\tts THICkness} & Calculates thickness from area integral\cr
{\tts WHATisit} & Identifies elements near a specified channel\cr
{\tts WIDth\_th} & Calculates thickness based on energy width\cr
{\tts INTSET} & Sets integration modes for {\tt THICKNESS} and {\tt INTEGRAL} commands \cr
{\tts DISPLAY} & Profiles the composition of the SIM sample structure \cr
\endautosizetable

The commands in this section are designed for analysis of simple elemental
RBS structures or for identification of features within a spectrum.  The
command generally use surface or mean energy scattering approximations.  For
more sophisticated analysis, the simulation capabilities described later
should be used.

Data for these commands is obtained from the {\tt ATOM3.DAT} file containing
atomic elemental information and stopping cross--section data.  Scattering
cross-sections are assumed to be purely Rutherford and resonances are not
treated by these commands.  Thickness, as measured by RBS, is an areal
density rather than a distance.  Conversion from the areal density of
atoms/cm$^2$ to thickness is performed using the standard atomic density of
the element.  Corrections for differing densities must be handled manually.

Many of these commands request input of marker positions from the cursor.  At
any point, the cursor input mode can be escaped by pressing the {\tt <esc>}
key and input will be taken from the keyboard instead.  From the graph, the
actual energy corresponding to the upper scale is obtained.  However, when
responding to a cursor request by typing, channel numbers should be specified.
These are translated to the equivalent energy using the {\tt CONV} parameters
of the ACTIVE buffer.  Inside macro files, it may be necessary to use {\tt
ENCURSOR NO} to completely disable the cursor modes and allow channel numbers
to passed through the command line.

<h2> Cursor </h2>

\syntax{CURsor}{Display cursor and return energy/channel values}

@@{CURSOR--RUMP}

The {\tt CURSOR} command displays a crosshair or hardware cursor on graphics
so equipped.  The coordinates of the crosshair are returned when any key or
button is pressed and are listed to the screen as channel number, energy and
yield (or counts if in raw mode).  If the key pressed was a {\tt 0} or {\tt
<sp>}, the cursor mode is retained for another value.  On any other key, the
cursor mode exist and the variables {\tts XCUR, YCUR, ECUR}, and {\tts CCUR}
are linked to the resulting channel, yield, energy (in MeV) and character
pressed.

The energy value returned by {\tt CURSOR} corresponds to the energy scale on
the graph.  However, the channel number shown corresponds to the channel in
the ACTIVE buffer which has the corresponding energy.  This may differ from
the lower channel scale on the axes if {\tt CONV} is different.

Ex: {\tt Channel:  338.4   Energy:  2.840 MeV   Yield:  28.1132 /uC/keV/msr}

\newpage
<h2> Info </h2>

\syntax{INFo <element>}{Print information about an element}

@@{INFO--RUMP}

The {\tt INFO} command prints a collection of information about the specified
element, including the isotopic makeup and scattering/stopping parameters.
The scattering geometry for cross section and stopping power evaluations is
taken from the ACTIVE buffer.  The first line after the parameters gives the
kinematic scattering factor and the channel and energy for scattering from an
element on the surface.  The matrix scattering height is the yield expected
from the pure element as a thick film on the surface.  The scattering cross
section is evaluated at the incident energy and a scattering angle of {\tt
180-phi}.  The stopping cross section {\tt [e]} includes corrections for the
incident and exit paths.  The stopping factor {\tt [S]} is determined using
the bulk elemental density.

\medskip\beginscreen
Your wish? \underbar{info si}

 -----------------------------------------------------
 Si  Z: 14  Mass:  28.09  Density:  0.4978E+23 at/cc ( 2.32 g/cc)

 Parameters: Energy  3.000 MeV  Theta  7.00     Phi   9.00
             keV/channel 4.950     keV(0)   1.600

 Surface Scattering:         0.5654 at  1.696 MeV (Channel: 342.3)
 Matrix scattering height:     9.05 Counts/uC/keV/msr
 Scattering cross section:     0.110 (1E-24 cm2/ster)
 Stopping Factors:  [e] =  76.3 (1E-15 eV-cm2)   [S] =  38.0 eV/A

Isotopes:  Mass:  29.97  Abundance: 0.03090
           Mass:  28.98  Abundance: 0.04700
           Mass:  27.98  Abundance: 0.92210

Ready if you are! 
\endscreen

<h2> Element </h2>

\syntax{ELement <element>}{Plots surface scattering energy}

@@{ELEMENT--RUMP}

The {\tt ELEMENT} command is used to identify on the graph the position where
an element on the surface would scatter.  The kinematic factor, energy and
channel for scattering are printed to the terminal.  If a graphic device is
active, a marker is drawn on the channel axis at the scattering energy and is
labeled with the element.  If an element name only is specified, the average
atomic weight is used.  Isotopes may be specified and will be marked
accordingly.

\medskip \screenwidth=\hsize \beginscreen
Feed me! \underbar{element si   el 28si}
 Si  Z=14  Mass= 28.086  K(He)=0.5654  Energy= 1.696 MeV  Channel= 342.321
 Si  Z=14  Mass= 27.977  K(He)=0.5641  Energy= 1.692 MeV  Channel= 341.550
At your service! 
\endscreen

<h2> Whatisit </h2>

\syntax{WHATisit <cursor\_entry>}{Identify elements near channel}

@@{WHATISIT--RUMP}

{\tt WHATISIT} is, in one sense, the inverse of the {\tt ELEMENT} command.
Whereas {\tt ELEMENT} is to mark a known elemental point, {\tt WHATISIT} is
used to mark 5 elements around a given energy position.  The five closest
elements to the marked point are displayed to the terminal and are drawn on
the graph.  These markers thus indicate at least the range of atomic numbers
corresponding to an unknown peak.  It is important to keep in mind that (1)
elements may not be on the surface and (2) that mass resolution is poor at
high Z and the actual element may not be listed at all.

\medskip\beginscreen
Ready if you are! \underbar{whatisit}

 Selected energy was  1.694 MeV
 Mg Mass  24.3   expected at  1.550 MeV (Channel  313)
 Al Mass  27.0   expected at  1.656 MeV (Channel  334)
 Si Mass  28.1   expected at  1.696 MeV (Channel  342)
 P  Mass  31.0   expected at  1.790 MeV (Channel  361)
 S  Mass  32.1   expected at  1.822 MeV (Channel  368)

Ready if you are! \underbar{whatisit}

 Selected energy was  2.769 MeV
 Ir Mass 192.2   expected at  2.762 MeV (Channel  558)
 Pt Mass 195.1   expected at  2.765 MeV (Channel  558)
 Au Mass 197.0   expected at  2.767 MeV (Channel  559)
 Hg Mass 200.6   expected at  2.771 MeV (Channel  560)
 Tl Mass 204.4   expected at  2.775 MeV (Channel  560)

At your service! 
\endscreen

<h2> Matrix </h2>

\syntax{MATrix <element>}{Plots an element's energy and height}

@@{MATRIX--RUMP}

A surface thick film of any pure element has a well defined yield height
determined by the stopping power, cross section and channel energy width. 
This height is referred to as the matrix height and is plotted and labeled on
the graph using the {\tt MATRIX} command.  The height and edge for the
specified element are plotted as a cross on the graph with the element name 
(or isotope) just above.  The height of a substrate edge, such as Si, can be
used to verify the normalization condition.  The scattering height and
position are also displayed on the terminal.  Matrix height is determined by
$$
     Yield =  { {\sigma * d\Omega * q *\Delta E_{chan}} \over {[\epsilon]}}
$$

\medskip\beginscreen
Feed me! \underbar{mat si}
 Si expected at  1.696 MeV ( 342.3) and height    9.052
At your service! 

\endscreen

<h2> CALIBRATE </h2>

\syntax{CALibrate <cursor\_pt> <cursor\_pt> <el> <el> <energy> <cursor\_pt>}{Calibrates ECONV and MEV for spectrum from specific peaks on spectrum}

The settings of the channel--to--energy conversion and the absolute energy of
the accelerator should be well known quantities but are often not controlled
well enough for an absolute setting.  The {\tt CALIBRATE} command determines
the values for {\tt CONV} and {\tt MEV} of the ACTIVE buffer from the
positions of two surface element peaks and one absolute energy point.
At Cornell, a $^{230}$Th alpha source is mounted near the detector to
provide the absolute energy calibration in all spectra.  If the beam energy is
accurately known, changes to the {\tt MEV} setting can be skipped by setting
the energy of the known marker to zero.  Only values for {\tt CONV} will be
modified in this case. 

When {\tt CALIBRATE} is executed, the cursor is enabled and two surface peaks
should be identified, either by a thin peak or the half height of a ledge. The
elements corresponding to these points are then entered (use the specific
isotope for better calibration especially if a low Z element is one of the
calibration points).  The absolute energy of a known marker (in keV) is then
requested.  The default value is the 4684 keV major peak from $^{230}$Th
decay.  If any positive non-zero marker energy is input, the cursor is again
displayed to obtain the position of the marker.  To improve accuracy, the
position of the marker should have previously been determined using the {\tt
REGION} and {\tt CURSOR} commands.  Escape out of cursor entry mode (using
{\tt <esc>}) and enter the absolute number from the keyboard.

The command changes the {\tt MEV} and {\tt CONV} values of the current buffer
to correspond to calculated values.  The accuracy should be tested using a
third element peak and the {\tt ELEMENT} command.  The following example
calibrates using buffer 3 and then transfers the calibration to real spectra
in buffers 1 and 2.

\medskip\beginscreen
Your wish?     \underbar{/* We first get the buffer and calibrate}
Up periscope!  \underbar{/* The channel for 230Th is escaped and typed}
Feed me!       \underbar{get 3 calibrate}
Identify two surface elements.
Element? (no change) \underbar{ni si}
Marker energy, keV (4684 for 230Th) ? \underbar{4684}
Marker position? \underbar{(pressed <esc> to abort cursor)}
Marker position? \underbar{930}
 Energy= 3.223 MeV    Conversion:   4.7825 keV/ch     3.24 keV(0)
Feed me! \underbar{for \%f in (1 2) do get \%f mev 3.223 econv 4.7825 3.24}
Beam me up Scottie! 
\endscreen

<h2> Integral/Thickness </h2>

\syntax{INTegral <cursor\_low> <cursor\_high>}{Integral of counts over a region}
\syntax{THICkness <cursor\_low> <cursor\_high> <element>}{Thickness calculation based on total counts in region}
\syntax{THICkness <cursor\_low> <cursor\_high> <element> <alpha\_val>}{}

@@{THICKNESS--RUMP} @@{INTEGRAL--RUMP}

The {\tt INTEGRAL} command determines the total integrated number of counts
between two markers in the ACTIVE buffer, either on a {\tt RAW} or {\tt
NORMALIZED} basis.  In the raw mode, the absolute counts in each channel are
summed.  In normalized mode, this sum is then scaled by the charge, detector
solid angle and correction factor ({\tt Q, OMEGA} and {\tt CORR}) to standard
conditions.  Two integrals, the ``gross'' and the ``net'', are reported. 
The gross integral is the complete sum while the net integral includes only
counts which lie above a straight line drawn between the region edges.   The
net integral thus performs the simplest of background subtractions.  However,
the {\tt BACKGROUND} command should be used to extract the peak first if it
lies on any substantial background signal.  

There are two ways in which the {\tt INTEGRAL} command can handle the channels
at either end of the integration region.  The {\tt INTSET} command is used
to set the mode.  In the default mode, the boundaries of the integration are
rounded to integers before summing the data (the rounding occurs toward the
center of the region).  {\tt INTSET} can switch the integration mode so that
the partial channels at the boundaries are also included in the integration.
See the {\tt INTSET} command for more details.

The {\tt THICKNESS} command first does an {\tt INTEGRAL} and then translates
the measured total yield into an equivalent thickness (atoms/cm$^2$) of an
element.   To first order, the total number of scattering centers (atoms) is
related directly to the integrated yield through the beam dose and detector
solid angle, $Y \approx {{d\sigma} \over {d\Omega}} * d\Omega * Q * N_t /
\cos \theta_{in}$.  Hence, the thickness determined by integration of a peak
is an absolute measurement assuming that $d\Omega$ and $Q$ are known
accurately.  Since these quantities
are usually known only to within about 5\%, the fudge factor {\tt CORR} also
enters into the equation.   Within these constraints, {\tt THICKNESS} provides
an absolute measurement of the number of atoms/cm$^2$ present.  This areal
density is presented as a thickness using the bulk elemental density. 
Thickness values are reported for both the gross and net integrals.

The {\tt INTSET} command also controls the manner in which the thickness is
calculated from the integral.  The default is to use the simple surface
scattering approximation.  Increasing accuracy is achieved by allowing {\tt
THICKNESS} to estimate the correction required because the scattering cross
section increases as the particle slows down in the layer.  The greatest
accuracy is obtained if the user provides the ``alpha'' energy loss ratio
factor (after setting the mode with {\tt INTSET Q}).  Even in the
compensated modes, the thicknesses are valid primarily for pure elemental
layers and {\tt SIM} should be used for more complex structures.  See the
{\tt INTSET} command for additional details on compensated thickness.

The variables {\tt GROSS\$}, {\tt NET\$}, and {\tt THICK\$} are linked by this
command. {\tt GROSS\$} and {\tt NET\$} are real numbers contain the
corresponding integrals.  {\tt THICK\$} is a four element array containing the
gross and net thicknesses in Angstroms (entries 1 and 2) and areal thickness
in atoms/cm$^2$ in entries 3 and 4.

\medskip\beginscreen \beginverb

Your wish? \underbar{thick ni}
Identify region with cursor.  Press any key.
     Discrete integration on Buffer  1,   Ni/NiSi/Si Annealed 90 min
 Region:  385.0 to  482.0 Gross:     5196.78 Net:     5164.94  (#/uC/msr)
 Ni surface approximation with micro-density of  8.90 g/cc yields
 (Gross)  0.1847E+19 Atoms/cm**2 ( 2023.4 Angstroms)
 ( Net )  0.1835E+19 Atoms/cm**2 ( 2011.0 Angstroms)

Yes dear? \underbar{intset interp thick ni intset round}
Identify region with cursor.  Press any key.
 Interpolated integration on Buffer  1,   Ni/NiSi/Si Annealed 90 min
 Region:  385.4 to  482.3 Gross:     5196.18 Net:     5158.00  (#/uC/msr)
 Ni surface approximation with micro-density of  8.90 g/cc yields
 (Gross)  0.1846E+19 Atoms/cm**2 ( 2023.2 Angstroms)
 ( Net )  0.1833E+19 Atoms/cm**2 ( 2008.3 Angstroms)

At your service! \underbar{intset estimate thick ni}
Identify region with cursor.  Press any key.
     Discrete integration on Buffer  1,   Ni/NiSi/Si Annealed 90 min
 Region:  385.0 to  482.0 Gross:     5196.78 Net:     5164.94  (#/uC/msr)
 Ni surface approximation with micro-density of  8.90 g/cc yields
 (Gross)  0.1847E+19 Atoms/cm**2 ( 2023.4 Angstroms)
 ( Net )  0.1835E+19 Atoms/cm**2 ( 2011.0 Angstroms)

Compensated calculation (refer to Chu et.al. page 65)
For fixed alpha =    1.13703      (surface approx. for Ni)
 (Gross)  0.1766E+19 Atoms/cm**2 ( 1934.8 Angstroms)
 ( Net )  0.1755E+19 Atoms/cm**2 ( 1923.3 Angstroms)

Your wish? 
\endverb \endscreen

<h2> Width\_thick </h2>

\syntax{WIDth\_th <cursor\_low> <cursor\_high> <element>}{Thickness
calculation based on peak width}

@@{WIDTH\_THICK--RUMP}

As an alternate to {\tt THICKNESS}, the width of a layer in energy can be
related to its thickness through the stopping factor [S].  The {\tt
WIDTH\_THICK} command determines the thickness based on a mean energy
approximation for the stopping factor.  (See Chu et al. for a discussion of
the mean energy approximation.)  {\tt WIDTH\_THICK} is appropriate for almost
pure elemental layers whose widths are substantially greater than the detector
resolution and form well defined flat top ledges.  The cursor is used to mark
the half--height points on either side of the layer.  The compensated stopping
factors and the calculated thickness are reported.  The accuracy of this 
calculation is constrained by the accuracy of the stopping power tables,
usually about 5-10\%.  The final thickness (\AA) and area thickness are linked
as the array {\tt THICK\$}.

\medskip\beginscreen
Feed me! \underbar{width\_th}
Place cursor at half-height points.
Element? (no change) \underbar{Ni}
 Stopping Factors:   [e] = 112.9 (1E-15 eV-cm2)  [S] = 103.0 eV/A
 Width:  254.5 keV    0.2254E+19 Atoms/cm**2     2470.3 Angstroms
Feed me! \underbar{eval thick\$(2)}
Value:   2.254E+18
Yes Master? 
\endscreen

<h2> INTSET </h2>

\syntax{INTSET [Round|Interp|Surf|Est|Query|?]}{Set mode for thickness calculations}

@@{INTSET--RUMP}

The {\tt INTSET} command sets the internal mode used for calculating the {\tt
INTEGRAL} or the {\tt THICKNESS} commands.  Integrals are evaluated in either
of two modes.  In the {\tt roundoff} mode, channels are rounded toward the
center of the integration region.  In the {\tt interpolate} mode, each channel
is treated as having a triangular basis extending one channel to either side
of center.   The edges of the region need not be integer channels and the
integral will include an appropriate fraction of the nearby edge channels. 
The default mode is {\tt rounding}.

Three modes may be set for determining the thickness of an element over the 
region, {\tt surface}, {\tt estimate} and {\tt query}.  The default is the
{\tt surface} approximation and no correction for the increase of the cross
section with depth is attempted.  The other modes {\tt estimate} and {\tt
query} are forms of the mean energy approximation.  The energy of the particle
at the time of scattering is estimated from the detected energy.  Energy loss
on the ingoing path is assumed to be a constant fraction of the total energy
loss and hence the scattering energy can be obtained from the overall energy
loss.  In terms of stopping powers, 
$$
E_{scat} = E_0-{\left ( {{dE}\over{dx}} \right )_{E_0} \over {cos \Theta_{in}}} *
          {{\Delta E} \over 
{\left [ K {\left ( {{dE}\over{dx}} \right )_{E_0} \over {\cos\theta_{in}}} +
           {\left ( {{dE}\over{dx}} \right )_{KE_0} \over {\cos\theta_{out}}} \right ]}}
$$
The parameter alpha is defined by the quantity
$$
 alpha = {{\cos\theta_{in}} \over {\cos\theta_{out}}} * 
           {{\left ( {{dE}\over{dx}} \right)_{KE_0}} \over 
           {\left ( {{dE}\over{dx}} \right)_{E_0}}}
$$
and is sufficient to perform the above calculations.  In the {\tt estimate}
mode, the thickness integrals are performed correcting for cross section
changes using this formulation of alpha.   For repeated measurements on a
given sample type, alpha can be calibrated accurately and the {\tt query}
mode allows for direct entry of a user value of alpha.  Given accurate user
values, {\tt THICKNESS} will correctly evaluate the areal density of elements
even in compounds.  See Chu et al. equations 3.21ff for additional details.

The ? command provides a listing and short description of these modes as well
as identifying which modes are currently active. 

\medskip\beginscreen
Feed me! \underbar{intset}
[Round|Interp|Surf|Est|Query|?] \underbar{?}

* current mode - Single letter command - Explanation of mode
* R - Integrals are rounded to nearest channel
  I - Integrals are interpolated between channels
* S - Thickness calculation based on surface approximation only
  E - Compensated thickness calculation with an estimated alpha
  Q - Compensated thickness calculation with query for alpha

At your service! \underbar{intset I intset E intset ?}

* current mode - Single letter command - Explanation of mode
  R - Integrals are rounded to nearest channel
* I - Integrals are interpolated between channels
  S - Thickness calculation based on surface approximation only
* E - Compensated thickness calculation with an estimated alpha
  Q - Compensated thickness calculation with query for alpha

At your service! 
\endscreen

<h2> Profile </h2>

\syntax{PROfile <element> <cursor\_region> <filename>}{Profile an element as concentration versus depth}

@@{PROFILE--RUMP}

The {\tt PROFILE} command (within RUMP --- not SIM) estimaes the concentration
versus depth for a single element corresponding over a given region in the
ACTIVE spectrum --- usually an isolate peak.  The stopping factors which
determine the non--linear conversion from energy to depth are calculated using
the current sample structure defined in the {\tt SIM} module.  A reasonable
sample structure {\it must be} defined prior to requesting {\tt PROFILE}.  For
example, to profile As as a dopant in Si, the sample structure need be defined
only as pure Si since the As will not strongly affect the stopping powers.

Concentrations are reported in atomic fraction and the depth is measured in
the only units appropriate to RBS, 10$^{15}$~at/cm$^2$.  {\tt PROFILE} makes
no attempt to deconvolute detector resolution, straggling, isotope effects, or
overlapping elements.  Hence, it is most useful for heavy, isolated impurity
peaks or monoisotopic elements.

The concentration versus depth is plotted if a graphics device is available.
After plotting, a filename is requested and the data may be written as ASCII
data pairs to the file.  Stored data can be subsequently massaged using
GENPLOT or other analysis programs.  If no filename (or a /) is given, the
data will not be stored.  The last command in the example below
might be used in a macro file for automatically profiling Ni over a given
range without saving the data.  The {\tt ENCURSOR} command allows the range
(in channel numbers) to be specified from the command line instead of by the
cursor (equivalent to escaping out of the cursor using {\tt <esc>}).

\medskip\beginscreen
Yes Master? \underbar{read {\bs}rump{\bs}examp}
File: c:{\bs}rump{\bs}examp.rbs
Id:   Ni/NiSi/Si Annealed 90 min 295C
Up periscope! \underbar{sim get {\bs}rump{\bs}examp return}
Beam me up Scottie! \underbar{profile}
Element to be profiled? (abort) \underbar{ni}
Use cursor to define region to be profiled.
File name (if any) of saved data: \underbar{niprof.dat}
Your wish? \underbar{encursor no profile ni 300 500 / encursor yes}
Hey, man, what next? 
\endscreen

Profile will fail if no simulation sample has been defined.  Also, the profile
will fail with the message {\tt Simulation goes to zero in region of interest}
if you request a profile over a region completely invalid for the element.
Invalid experimental parameters will also generate an error.  Warnings
are given if the profile is subject to some isotope effect or if absurd
answers would be expected because of the isotopic distribution (profiling B
for example). 

\hl1{Other Commands}

The following system utility commands are identical between GENPLOT and RUMP 
and are described in other sections of the manual as specified in the table
below.

\medskip\autowraptable{.25\hsize}{.75\hsize}*
\multicol{2}{l}{\bf Section J -- Function Evaluation}\cr
{\tts ALLOCate} & Allocate a variable \cr
{\tts DEALLocate} & Deallocate a variable \cr
{\tts DEClare} & Allocate and set a STRING variable \cr
{\tts DEFine} & Define a user function \cr
{\tts EValuate} & Evaluate a general math expression \cr
{\tts LET} & Set a variable to a constant or expression \cr
{\tts LISTVar} & List all the defined variables \cr
{\tts SAVEVars} & Save variables to disk file \cr
{\tts SETVar} & Allocate and set a real variable \cr
\cr
\multicol{2}{l}{\bf Section K -- Command Macros}\cr
{\tts BASE\_MACRO} & Specify a default macro directory \cr
{\tts ECHO} & Toggle the command display while a macro is running \cr
{\tts FOR} & Repeat command loop \cr
{\tts MACRO} & Collect commands in a macro file \cr
{\tts GOTO} & Continue execution at the given label \cr
{\tts IF} & Conditional execution of commands \cr
{\tts QUERY} & Prompt the user for input from the macro \cr
{\tts SAVE\_MACRO} & Save the commands from MACRO in a file \cr
{\tts SLEEP} & Delay for a given length of time \cr
{\tts WAIT} & Wait for user action before continuing \cr
{\tts XEQ} & Execute a macro file \cr
{\tts CMDLIN} & Execute the next token as a series of commands \cr
{\tts \&GETARG} & Get the next argument from the macro command line \cr
{\tts \&QUERY} & Get the next token from the user \cr
{\tts \&YESNO} & Get a yes/no answer from the user \cr
{\tts \&ENCODE} & Real number to string encoding \cr
\endautowraptable

\medskip\autowraptable{.25\hsize}{.75\hsize}*
\multicol{2}{l}{\bf Section L --- Operating System} \cr
{\tts CD} & Change directory \cr
{\tts DIr} & Display directory \cr
{\tts <disk>:} & Change default disk \cr
{\tts DOS} & Send a command to the system to be executed \cr
{\tts MEMory} & Display unused RAM space \cr
{\tts PAUSE} & Temporarily pause GENPLOT and bring in the system \cr
{\tts TYpe} & Type a file on the screen \cr
{\tts WHEre} & Gives the pathname of the current directory \cr
\cr
\multicol{2}{l}{\bf Section N --- Editing}\cr
{\tts EMacs} & Built in editor\cr
\endautowraptable

\closesyntax
\bye

\input rpsim
\input rppert

%------------------------------------------------------------

\hl1{Starting RUMP and Initialization}

<h2> Starting RUMP </h2>

\syntax{RUMP [ini\_file]}{Start RUMP with optional initialization file}
\syntax{PERT [ini\_file]}{Start RUMP and PERT with optional initialization file}

RUMP and PERT may be started directly from the {\tt .EXE} files distributed in
the {\tt RUMP} directory.  RUMP and PERT differ only in that RUMP has all
routines associated with the PERT commands nulled out to save space.  Hence, 
RUMP is a complete subset of PERT.  Both RUMP and PERT require a full 640 kB
configuration and will not run if a substantial number of terminate and stay
resident routines have been loaded (TSR's).  

Because of the large memory requirements of RUMP and PERT, they are normally
executed from the batch files {\tt RUMP.BAT} or {\tt PERT.BAT}.  These batch
files define an environment variable {\tt HDATA.CHR} which causes only 4
character sets to be loaded into the graphic plotting routines.  The batch
files point RUMP initialization to {\tt RUMP.INI} and PERT initialization to
{\tt PERT.INI} or as specified on the calling command line.


<h2> Initialization </h2>

The default directory for RUMP internal data files and the initialization file
is {\tt c:{\bs}rump}.  RUMP will search for atomic stopping power data files,
resonance cross section tables and the initialization files in this
directory if not found elsewhere.  The search order is (1) check for an
environment variable with the same name as the requested file and use the
value of the environment variable as the actual filename, (2) check for file
in the current directory and (3) check for file in the default directory.

RUMP automatically loads the {\tt ATOM3.DAT} file during initialization.  This
file must contain the atomic stoppping power data.  Next, RUMP searches for 
the default initialization file {\tt RUMP.INI} or the initialization file
specified on the command line.  If an initialization file is found, it will be
executed as a macro file after other initialization.  The file is always
executed with {\tt ECHO OFF}.  Initialization files should configure RUMP with
the appropriate facility dependent parameters.

<h2> RUMP Data Files </h2>

RUMP requires only a single data file, the atomic stopping power and isotope
information file.  The stopping power tables are initially loaded from
{\tt atom3.dat} in the default directory.  These tables, as distributed, are
configured for alpha particle energies between 0.35 and 3.5 MeV.  See the
{\tt LOAD} command and the {\tt QSTOPP} and {\tt STOPP} programs for
generating alternate stopping power tables.  Resonance cross section files are
also, by default, located in the {\tt c:{\bs}rump} directory and loaded as
requested by the {\tt LOAD} command.

Many other data files are part of the underlying subsystem and are loaded from
the {\tt c:{\bs}genplot} directory.  RUMP only files which are located in 
{\tt c:{\bs}genplot} are include {\tt RUMP\_RDR.OVL} and {\tt RUMP\_DYN.OVL}.
These files contain additional executable code for {\tt RUMP.EXE} and {\tt
PERT.EXE}. 

<h2> VAX Specifics </h2>
%
%@@{VAX--RUMP}
%
%A filename on the VAX is the title of a clump of information.  In particular,
%your spectrum with parameters is stored as a file. The name of a file must be
%{\tt TITLE.EXT}, where {\tt TITLE} is 1--9 characters, letters or numbers.
%{\tt EXT} is 0--3 characters, letters or numbers. No distinction is made here
%between upper and lower case letters.  Example: {\tt ALCU1G.RBS} would be a
%glancing angle RBS spectrum.  RUMP will supply the {\tt .RBS}, but VMS
%commands do not.
%
%The version number of a file is shown by a {\tt ;N} at the end of the file
%name.  You can usually ignore this -- all commands refer to the latest
%version.  For {\tt DELETE}, however, you must specify which version (use {\tt
%;*} for all of them).
%
%The wildcard characters for use with {\tt GET} are `*' and `?'.  `*' will
%match any string of characters, and `?' will match any single character.
%
%\hl3{Spawn}
%
%Spawn a subprocess to execute VMS commands.
%
%{\tt SPAWN} can be particularly useful in editing macro files.  Simply say
%{\tt spawn~edit~filename}, and when you exit the editor you can {\tt
%xeq~filename} to try it out.  This is convenient when you are working on a
%macro to do some complex plotting.
%
%\hl3{Initialization File}
%
%RUMP will search for the file {\tt SYS\$LOGIN:RUMP.INI} when it is started up.
%If found, this file will be run as a macro in {\tt QUIET} mode.
%%\bye
%you are working on a
%macro to do some complex plotting.
%
%\hl3{Initialization File}
%
%RUMP searches for the file {\tt SYS\$LOGIN:RUMP.INI} when it is started up.
%If found, this file will be run as a macro in {\tt QUIET} mode.

\input rpfile
\closesyntax
</html>
